\section{模型评价与推广}
\subsection{模型优点}
模型对四问采用同一数据口径：瞬时 GPU 容量按任务重叠计算，小时电量按持续时间分摊，避免平均负荷掩盖峰值。Q1 严格按训练、验证、闭环测试排序，并保留季节朴素基线；Q2 将增量候选评价与 50000 个任务的全量审计分开；Q3 强制保留 no-action、对轨迹去重；Q4 在每轮储能更新后重建任务邻域。这些设计分别控制了数据泄漏、搜索成本、基线缺失和“伪联合优化”风险。

五个任务方案、七个储能策略和 17 个联合情景均有逐行结构化结果，硬约束残差通过；两次隔离运行的结果哈希及 Q2--Q4 语义一致。因此结论可复算，并达到当前证据边界下的 scenario-feasible。

\subsection{模型不足}
Q2 每个方向只检查 10000 个单任务移动，虽单方案覆盖超过 3\%、五方案合计超过 16\%，仍未穷尽任务—时空组合；Q4 每轮只检查 600 个候选并最多运行 8 轮，二者均没有全局最优证书。Q3 的 21 状态 SOC 离散也只给出网格内候选。储能模型未计退化、自放电和备用容量，部分策略等效循环超过 1500，故账面收益不能直接解释为可部署收益。负成本主要来自题面边界内售电收入高于购电支出，也不表示全生命周期成本为负。

通信侧仅使用静态单向时延，未纳入带宽、拥塞和迁移能耗；最大等待仍可达到 219 小时，说明综合服务指标不能替代任务级尾部检查。统一 21 状态基准下碳排放为 0 只说明题面边界和当前候选允许该结果，不代表生命周期零碳。当前 benchmark 又无独立 Oracle，约束通过和双运行只能证明内部一致与可重复，不能证明外部真实最优。

\subsection{推广方向}
后续应优先扩大 Q2/Q4 邻域或引入混合整数规划上界，并沿用同一全量审计；Q3 可加入每 MWh 吞吐退化成本、年度循环上限并做网格加密。在线场景可采用滚动时域：逐小时更新任务、新能源和价格预测，冻结已开工任务并传递绝对 SOC，这与模型预测控制一致\cite{rawlings2017model}。多目标部分可用 $\varepsilon$-约束法扩充非支配候选\cite{miettinen1999nonlinear}，但仍须区分“候选内非支配”和“已证 Pareto”。
